\documentclass[11pt]{article}
\usepackage{amsmath,amssymb,amsthm,fullpage}
\newtheorem{theorem}{Theorem}
\newtheorem{lemma}{Lemma}
\title{A Shared Four-Term Pool and Twenty Tree Queries}
\author{}
\date{November 24, 2026}
\begin{document}
\maketitle

\paragraph{Verification status and scope.}
The results here are complete written proofs about masked exterior Cauchy
kernels. The companion Python checker uses exact rational arithmetic for
the finite certificate and the limiting inequalities. These proofs are
not Lean checked. They neither prove a superconstant saving nor establish
an unrestricted lower bound for HeadComplexity. No model definition is
changed. The main new point is stronger than eight-query shattering with
four terms: a single four-term pool, after its four full-sum signs have
been prescribed, realizes every proper-descendant pattern by permutation.
This gives a target-adaptive eight-term construction with twenty signs.

\section{The shared pool}
Put $\eta=1/1000$. The four full-sum query times, two pair query times,
and singleton query time are respectively
\[
 t_j=1+\eta u_j,\quad u=(-3,-1,1,3),\qquad
 t_-=1-2\eta,\quad t_+=1+2\eta,\qquad t_0=1.
\]
For four terms $r_h/(t-q_h)$ and a permutation $\pi$ of
$\{0,1,2,3\}$, the four proper-descendant scores are
\[
 \sum_{i\in\{1,3\}}\frac{r_{\pi(i)}}{t_--q_{\pi(i)}},\quad
 \sum_{i\in\{1,3\}}\frac{r_{\pi(i)}}{t_+-q_{\pi(i)}},\quad
 \frac{r_{\pi(2)}}{t_0-q_{\pi(2)}},\quad
 \frac{r_{\pi(3)}}{t_0-q_{\pi(3)}}.
 \tag{1}\label{eq:internal}
\]

\begin{theorem}\label{thm:pool}
There exists $\epsilon_0>0$ such that for every
$0<\epsilon<\epsilon_0$ and every prescribed full-sum sign vector
$s\in\{-1,1\}^4$, the fixed poles
\[
 q_h=\epsilon a_h+\epsilon^2 b_h,\qquad
 a=(-3,-1,1,3),\quad b=(0,16,0,0)
 \tag{2}\label{eq:poles}
\]
admit nonzero real residues $r_h$ with these two properties:
\begin{enumerate}
\item $\sum_h r_h/(t_j-q_h)$ has sign $s_j$ for every $j$;
\item as $\pi$ varies, the four scores in \eqref{eq:internal}
realize all $16$ strict sign vectors.
\end{enumerate}
All denominators are positive. The same $\epsilon_0$ works for every
$s$.
\end{theorem}

\begin{proof}
Write $P_\epsilon(x)=\prod_h(x-q_h)$ and choose a cubic polynomial
$N(x)=\prod_{\ell=1}^3(x-\rho_\ell)$, with the possible additional
factor $-1$. Set
\[
 r_h=\frac{N(q_h)}{P_\epsilon'(q_h)}.
\]
The partial-fraction identity gives
\[
 \sum_h\frac{r_h}{t-q_h}=\frac{N(t)}{P_\epsilon(t)}.
 \tag{3}\label{eq:partial}
\]
Consequently the full-sum signs are exactly the signs of $N$ at the
four query times.

Here is a finite list of eight root triples. An entry $v$ in the table
means $\rho=1+\eta v/10$. The last column records the positive-sign
bit pattern of $N(t_j)$, with the first query as the least significant
bit.
\[
\begin{array}{c|r r r|r}
 &v_1&v_2&v_3&\text{pattern}\\ \hline
1&-34&-1&44&3\\
2&17&-15&-12&8\\
3&-7&-24&45&2\\
4&26&-27&-2&10\\
5&-52&23&14&15\\
6&43&-36&-16&1\\
7&-42&19&6&11\\
8&33&12&-27&6
\end{array}
\]
Together with the negatives of these polynomials, these give all sixteen
full-sum patterns. Direct rational substitution gives, for every row,
\[
 t_-< t_*:=\frac{1}{\rho_1^{-1}+\rho_2^{-1}+\rho_3^{-1}-2}<t_+.
 \tag{4}\label{eq:zero}
\]
These finitely many elementary inequalities are independently checked
by the companion script.

We verify that each chosen polynomial has the required permutation
property for all sufficiently small $\epsilon$. Put $N_0=N(0)$,
$N_1=N'(0)$, and
\[
 v_h=\frac1{\prod_{j\ne h}(a_h-a_j)}
       =\frac1{48}(-1,3,-3,1)_h,
 \qquad
 L_h=\sum_{j\ne h}\frac{b_h-b_j}{a_h-a_j}=(8,-4,-8,-4)_h.
\]
Taylor expansion of a finite rational expression, with all displayed
limiting denominators nonzero, yields
\[
 r_h=\epsilon^{-3}v_hN_0+
       \epsilon^{-2}v_h(a_hN_1-N_0L_h)+O(\epsilon^{-1}).
 \tag{5}\label{eq:residue}
\]
The residues therefore alternate in sign for small $\epsilon$.
The four pairs having opposite residue signs are
$\{0,1\},\{0,3\},\{1,2\},\{2,3\}$.
For the first and last of these, multiplication of the partial score
by $\epsilon^3$ gives the nonzero limits
\[
 \frac{N_0}{24t},\qquad -\frac{N_0}{24t}.
\]
They consequently give the two constant pair-sign patterns at
$(t_-,t_+)$.

For the outer pair $S=\{0,3\}$ one has
\[
 \sum_{h\in S}v_h=0,\quad
 \sum_{h\in S}a_hv_h=1/8,\quad
 \sum_{h\in S}v_hL_h=-1/4.
\]
Using \eqref{eq:residue} and expanding $(t-q_h)^{-1}$ gives
\[
 \epsilon^2\sum_{h\in S}\frac{r_h}{t-q_h}
 \longrightarrow
 \frac{N_0+(N_1+2N_0)t}{8t^2}.
 \tag{6}\label{eq:limit}
\]
For the inner pair $\{1,2\}$ the limit is the negative of this
expression. Since $N_1/N_0=-\sum_\ell\rho_\ell^{-1}$,
the numerator in \eqref{eq:limit} vanishes precisely at $t_*$.
By \eqref{eq:zero} the two symmetric pairs give the two alternating
pair-sign patterns for small $\epsilon$.

Every one of the four possible pair-sign patterns is thus supplied by
a pair containing one residue of each sign. Its complementary pair
also contains one residue of each sign. After choosing the pair occupying
positions $1,3$, its internal order sets the sign at position $3$
arbitrarily. Independently ordering the complementary pair sets the
sign at position $2$ arbitrarily. This proves all sixteen patterns
in \eqref{eq:internal}. Positivity and nonvanishing hold for sufficiently
small $\epsilon$. Finally take the minimum of the finitely many positive
thresholds for the eight polynomials and their negatives.
\end{proof}

\section{Adding full-sum signs by paired cancellation}
\begin{lemma}\label{lem:jets}
Let $p_1,\ldots,p_M$ be distinct real numbers below $2M$ distinct
query times $x_1,\ldots,x_{2M}$, and let $c_h\ne0$.
Start with the $2M$ terms
\[
 \frac{c_h}{t-p_h},\qquad -\frac{c_h}{t-p_h},\quad 1\le h\le M.
\]
Any prescribed strict signs at all $2M$ full-sum query times can be
obtained by an arbitrarily small perturbation of just the poles and
residues of the negative copies. Any finite collection of strict
partial-sum signs at the initial configuration is preserved by a
sufficiently small such perturbation.
\end{lemma}
\begin{proof}
Perturb the negative copies to
$(-c_h+\delta\alpha_h)/(t-p_h-\delta\beta_h)$.
The derivative at $\delta=0$ of the full sum is
\[
 \sum_h\left(\frac{\alpha_h}{t-p_h}
              -\frac{c_h\beta_h}{(t-p_h)^2}\right).
\]
The evaluation matrix of these $2M$ functions at the $2M$ query
times is invertible. Indeed, multiply a purported vanishing linear
combination by $\prod_h(t-p_h)^2$. The result is a polynomial of
degree at most $2M-1$ with $2M$ distinct roots, hence zero. Uniqueness
of the principal parts at the distinct poles then makes every
coefficient zero. Since $c_h\ne0$, arbitrary derivative values at
the query points are available. Choose them equal to the desired
signs. For sufficiently small positive $\delta$, the full sums have
those strict signs. The additional strict signs and all denominator
inequalities persist by continuity on a finite set.
\end{proof}

\begin{theorem}\label{thm:twenty}
For each sufficiently small common $\epsilon$, there are twenty query times
fixed across all target sign vectors, and fixed masks on eight terms, such
that every twenty-bit sign vector is realized by masked exterior Cauchy
sums with label-dependent poles and residues. The pole magnitudes can
be arbitrarily small, and all query times lie in $[1-7\eta,1+7\eta]$
after subtracting a common shift tending to zero.
\end{theorem}
\begin{proof}
Use eight full-sum query times $1+\eta u$ with
$u=-7,-5,-3,-1,1,3,5,7$. Divide the eight terms into a left group
and a right group of four. Query the right group at the four full-sum
times of Theorem~\ref{thm:pool}. In each of the two groups query the
four proper descendants in \eqref{eq:internal}. This gives
$8+4+4+4=20$ fixed queries.

Choose the shared pool from Theorem~\ref{thm:pool} using the four
prescribed right-group full-sum signs. In the right group, permute that
pool to obtain its four prescribed descendant signs. In the left group,
put a negative copy of the same pool and permute it to obtain its four
prescribed descendant signs. The universal permutation property makes
these two choices independent. At this stage twelve required signs
are strict, and the full eight-term sum is identically zero.
Lemma~\ref{lem:jets} supplies all eight remaining full-sum signs while
preserving those twelve signs.

For positive small poles, add $4\epsilon$ simultaneously to all poles
and query times before applying the lemma. The smallest unperturbed
pole is then $\epsilon$. Choose the perturbation small enough to keep
all poles positive and of size $O(\epsilon)$. All query denominators
remain positive. As $\epsilon$ can be arbitrarily small, so can the
final pole magnitudes.
\end{proof}

\section{An original-model axis-zero consequence}
\begin{theorem}\label{thm:compiler}
For every $r\geq1$ and every sign table on an $r\times20$ grid,
there is an axis-zero compiler with sixteen affine-ratio heads. Every
denominator has strictly negative slopes and is positive on the entire
Boolean $(r+20)$-cube.
\end{theorem}
\begin{proof}
Choose a common
\[
 0<\epsilon<\min\{\epsilon_0,10^{-6},1/(10000r)\}.
\]
Use Theorem~\ref{thm:twenty} independently for each of the $r$ rows,
with the same fixed twenty times and masks. In every row take the final
pole perturbations smaller than $\epsilon/2$, and preserve the strict
signs as in the proof. Add $4\epsilon$ to the centered poles and query
times as there, then multiply all poles and times by
$\alpha=1000/1008$. Multiply all residues by $\alpha$ too, preserving
the scores exactly. The resulting row poles $p_{hi}$ and common query
times $t_j$ satisfy
\[
 0<p_{hi}<8\epsilon,\qquad .985<t_j<.9991.
\]
Let $a_{hi}$ be the resulting residues, and let $M_{hj}\in\{0,1\}$
be the fixed masks. Put $\zeta=10^{-4}$ and define, for each of the
eight pole indices $h$, two affine denominators and one numerator:
\[
\begin{split}
 B_h^+(z,w)&=1+\zeta-\sum_i p_{hi}z_i
               -\sum_j\bigl(\zeta+M_{hj}(1-t_j)\bigr)w_j,\\
 B_h^-(z,w)&=1+\zeta-\sum_i p_{hi}z_i-\sum_j\zeta w_j,\\
 A_h(z,w)&=\sum_i(1-p_{hi})a_{hi}z_i.
\end{split}
\]
Every slope is strictly negative. Moreover
\[
 \sum_i p_{hi}+20\zeta+\sum_j M_{hj}(1-t_j)
 <.0008+.002+.3<1+\zeta.
\]
Thus both denominators are positive on the entire cube. The sixteen-head
score is
\[
 F(z,w)=\sum_{h=1}^8\left(\frac{A_h(z,w)}{B_h^+(z,w)}
                         -\frac{A_h(z,w)}{B_h^-(z,w)}\right).
\]
It vanishes whenever $z=0$, because the numerators vanish, and whenever
$w=0$, because the two denominators in each pair agree. At a grid point,
a direct subtraction gives
\[
 F(e_i,e_j)=(1-t_j)\sum_{h:M_{hj}=1}\frac{a_{hi}}{t_j-p_{hi}}.
\]
The factor $1-t_j$ is positive, proving the prescribed strict sign.
This exact paired identity and its full-cube axis-zero admissibility
are independently Lean checked as
\texttt{HeadComplexity.AxisZeroGluing.exists\_maskedCauchy\_reverseAxisData}
in \texttt{Results/ReverseMaskedCauchy.lean}. The analytic shared-pool
existence theorem used here has not been formalized, so this complete
compiler corollary remains a written proof.
\end{proof}

\paragraph{Why this is not an induction.}
Theorem~\ref{thm:pool} is a uniform four-term permutation-pool theorem.
Theorem~\ref{thm:twenty} gives a target-adaptive eight-term representation;
it does not produce one eight-term pool whose permutations realize all
$4096$ proper-descendant patterns. That stronger invariant, or an
adequate target-adaptive replacement, would be needed to repeat the
argument. Exhausting all $40320$ permutations of the exact positive and
negative doubles of the eight displayed four-term pools gives only
$800$, $884$, or $1172$ strict proper-descendant patterns. This finite
failure does not exclude different eight-term pools or other mechanisms.
The missing unbounded construction is not assumed here.

\end{document}
